%%%%%%%%%%%%%%%%%%%%%%% file template.tex %%%%%%%%%%%%%%%%%%%%%%%%%
%
% This is a general template file for the LaTeX package SVJour3
% for Springer journals.          Springer Heidelberg 2010/09/16
%
% Copy it to a new file with a new name and use it as the basis
% for your article. Delete % signs as needed.
%
% This template includes a few options for different layouts and
% content for various journals. Please consult a previous issue of
% your journal as needed.
%
%%%%%%%%%%%%%%%%%%%%%%%%%%%%%%%%%%%%%%%%%%%%%%%%%%%%%%%%%%%%%%%%%%%
%
% First comes an example EPS file -- just ignore it and
% proceed on the \documentclass line
% your LaTeX will extract the file if required
\begin{filecontents*}{example.eps}
%!PS-Adobe-3.0 EPSF-3.0
%%BoundingBox: 19 19 221 221
%%CreationDate: Mon Sep 29 1997
%%Creator: programmed by hand (JK)
%%EndComments
gsave
newpath
  20 20 moveto
  20 220 lineto
  220 220 lineto
  220 20 lineto
closepath
2 setlinewidth
gsave
  .4 setgray fill
grestore
stroke
grestore
\end{filecontents*}
%
\RequirePackage{fix-cm}
%
%\documentclass{svjour3}                     % onecolumn (standard format)
%\documentclass[smallcondensed]{svjour3}     % onecolumn (ditto)
\documentclass[smallextended]{svjour3}       % onecolumn (second format)
%\documentclass[twocolumn]{svjour3}          % twocolumn
%
\smartqed  % flush right qed marks, e.g. at end of proof
%
\usepackage{graphicx}
%
% \usepackage{mathptmx}      % use Times fonts if available on your TeX system
%
% insert here the call for the packages your document requires
%\usepackage{latexsym}
% etc.
%
% please place your own definitions here and don't use \def but
% \newcommand{}{}
%
% Insert the name of "your journal" with
% \journalname{myjournal}
%
\begin{document}

\title{Developing a parsimonius predictor for binary traits in sugar
  beet (\emph{Beta vulgaris}) \thanks{Filippo Biscarini and Nelson Nazzicari contributed
    equally to the work.}
}
%\subtitle{Do you have a subtitle?\\ If so, write it here}

\titlerunning{Parsimonious predictor for binary traits in sugar beet}        % if too long for running head

\author{Filson Nazzarini         \and
        Simone Marini \and
	Piergiorgio Stevanato \and
	Nelppo Biscicari %etc.
}

%\authorrunning{Short form of author list} % if too long for running head

\institute{F. Biscarini \at
              Fondazione Parco Tecnologico Padano \\
              Tel.: +123-45-678910\\
              \email{filippo.biscarini@tecnoparco.org}           %  \\
%             \emph{Present address:} of F. Author  %  if needed
 \and
 S. Marini \at
    second address
	\and
	N. Nazzicari \at
    Fondazione Parco Tecnologico Padano \\
    Tel.: +123-45-678910\\
    \email{nelson.nazzicari@tecnoparco.org}           %  \\
	\and
	P. Stevanato \at
	address
}


\date{Received: 05 August 2014 / Accepted:}
% The correct dates will be entered by the editor

\maketitle

\begin{abstract}
Insert your abstract here. Include keywords, PACS and mathematical
subject classification numbers as needed.
\keywords{binary traits \and genomic predictions \and parsimonious predictor \and sugar beet}
% \PACS{PACS code1 \and PACS code2 \and more}
% \subclass{MSC code1 \and MSC code2 \and more}
\end{abstract}

\section{Introduction}
\label{intro}
The primary goal of breeding schemes in farm animals and crops is
generally to increase the agricultural output. Production traits are
typically quantitative continuous variables (e.g. milk
yield in dairy cattle, or per hectare yield in maize and rice).
Many traits of importance in plant and animal breeding follow nonetheless
a discrete categorical distribution, both ordered (e.g. calving ease in
cattle, grain texture in rice) and unordered
(e.g. grain pigmentation in rice, coat colour in cattle). A special case
is that of binomial traits, which can take up only two different values,
like disease resistance/susceptibility or presence/absence of a
morphological characteristic. 
Annual bolting (flowering) behaviour and root vigor are examples of binomial traits of agronomic
importance in sugar beet (\emph{Beta vulgaris}). 

Advances in biotechnology and genomics, and the advent of high-density
molecular markers (especially sinlge-nucleotide polymorphisms, SNP)
genotyping have led to the emergence of molecular breeding
\cite{moose2008molecular}.

One exciting application of molecular breeding are genomic predictions
(\cite{meuwissen2001prediction}): something on genomic predictions in
general. The recent publication of the reference genome for \emph{Beta
vulgaris} genome \cite{dohm2013genome} facilitated molecular breeding
also in this crop species. Genomic predictions in sugar beet already
done both for continuous (\cite{hofheinz2012genome,wurschum2013genomic}) and binary
(\cite{biscarini2014genome}) traits.

The concept of parsimony: when many possible predictors are available, it is
useful to select a subset to limit analysis cost and time. Moreover: use
the minimun necessary information set, occam razor (\cite{chaitin2006limits}), and so forth.
A model need to be simpler than the data the it fits/explains (e.g. knn with k=1)

Given two models that fit the data, the simplest has to be chosen (Occam's razor)


As the technology advances, and available predictors grow, not only the 
prediction precision becomes important, but also the actual cost must
be considered.

Sugar beets in particular: we work on root vigor \cite{biscarini2014genome}.

In this paper we propose statistical methods to highlight and select the 
most useful predictors given a set. We started on real world data and validated
our approach on a XXX dataset. We found that it is possible to strongly reduce
the dimension of the predictors set and still achieve high performance.

\section{Material and methods}
\label{sec:1}
\subsection{Plant material and SNP genotypes}
\label{sec:data}
Root vigor. Available data. SNP technology used, imputation. \\
Copypaste from other articles. Dataset description.
Text with citations \cite{stevanato2013high} and
\cite{saccomani2009molecular}.

\subsection{Predictor development procedure}
\label{sec:overview}
A two-step approach was adopted for the construction of a parsimonious
predictor for root vigor.

- a ranker to rank the various available predictors (SNPs in our case). We
  used the BOSS algorithm
- this is an iterative step. we progressively reduced the predictors set, 
  taking away the laest useful predictor and applying to the resulting
  subset a ridge logistic regression apprach. Thus, we obtained as many
  performances estimation as the number of original predictors.
\subsubsection{Rank of predictors}
\label{par:boss}
This explain the BOSS algorithm \cite{russu2012stochastic}

\subsubsection{Selection of predictors and classification method}
\label{par:predictor_selection}
We take one predictor out at each iteration
You put the model formula for ridge logistic regression

\subsubsection{Predictive ability}
\label{par:estimating_error}
Cross validation: how many times, what fractions. 
Explanation of error rate and other parameters (ROC?)

\subsection{Comparison with another method to rank predictors}
\label{sec:other_ranker}
Another ranker: why use one, and its description.
P value and SNP effect (as it is done in GWAS)

SNP variance \cite{gianola2009additive}

\subsection{Software}
\label{sec:software}
R, weka, perl.

\section{Results}
\label{sec:results}
Possible charts:
- Precision as a function of the number of predictors.
- Breakdown of two types of error.

If possible: probability of assignment as a function of predictors, maybe
with ROC curve? Maybe in discussion?

Plot AUC vs \# predictors?

Something on the ranking of predictors (variable importance vs SNP effects)?

Plot coefficients from BOSS?

% For one-column wide figures use
%\begin{figure}
% Use the relevant command to insert your figure file.
% For example, with the graphicx package use
%  \includegraphics{example.eps}
% figure caption is below the figure
%\caption{Please write your figure caption here}
%\label{fig:1}       % Give a unique label
%\end{figure}

\section{Discussion}
\label{sec:discussion}
General overview
why error rates are not evenly distributed?
Reminder: it works very well because of LD and H2

\subsection{SNP effects}
Manhattan plot with BOSS weights and weights from the other articles,
somehow compared (same chart? two charts? only ten best?).

Do the peaks make sense from the biological perspective?

Variance of SNPs vs genetic variance: $\rightarrow$ missing
heritability? (cite Brachi 2011, Manolio 2009?).

BOSS probability: 1 big peak + smaller peaks. Compare against SNP
density? Maybe the big peak corresponds to a physically isolated SNP,
whereas smaller peaks correspond to a cluster of SNPs in LD which
individually account for a smaller part of the variatino, but together
play an important predictive role. 

\subsection{Relative performance of rankers}
why using Pvalues and not other standard rankers (e.g. backward stepwise
selection)? Because of the specific nature of the problem

\subsection{Genotyping strategies and applications to breeding}
genotyping strategies: 
Costs, possible technologies (gbs, snp chip, macroarrays), implications

applications to breeding:
why is it important root vigor early detection. Other binomial traits (e.g.
disease resistance) May be applied to bolting (another trait which
exhibits binomial distribution), which has been shown to be controlled
by multiple genes and influenced by environmental factors
(\cite{salah2012genetic}).

sugar beet: $30\%$ of world's sugar production (cite Dohm? FAO?). Root
vigor linked to yield.

Sugar beet: sugar + energy (citation?)

Other binomial traits: resistance to viral and fungal diseases, bolting
(cite Dohm? Someone else?)

Breeding has shaped the genome of sugar beet (comparison with \emph{Beta
  maritima}, \cite{dohm2013genome}).

Extensions to multinomial traits? Examples?

% For two-column wide figures use
%\begin{figure*}
% Use the relevant command to insert your figure file.
% For example, with the graphicx package use
%  \includegraphics[width=0.75\textwidth]{example.eps}
% figure caption is below the figure
%\caption{Please write your figure caption here}
%\label{fig:2}       % Give a unique label
%\end{figure*}
%
% For tables use
%\begin{table}
% table caption is above the table
%\caption{Please write your table caption here}
%\label{tab:1}       % Give a unique label
% For LaTeX tables use
%\begin{tabular}{lll}
%\hline\noalign{\smallskip}
%first & second & third  \\
%\noalign{\smallskip}\hline\noalign{\smallskip}
%number & number & number \\
%number & number & number \\
%\noalign{\smallskip}\hline
%\end{tabular}
%\end{table}

\section{Conclusions}
\label{sec:conclusions}

Concluding remarks

\begin{acknowledgements}
Do we need to ack somebody? (projects?)
This research was financially supported by the Marie Curie European
Reintegration Grant ``NEUTRADAPT''.
\end{acknowledgements}

% BibTeX users please use one of
%\bibliographystyle{spbasic}      % basic style, author-year citations
\bibliographystyle{spmpsci}      % mathematics and physical sciences
%\bibliographystyle{spphys}       % APS-like style for physics
\bibliography{parsimonious.bib}   % name your BibTeX data base

% Non-BibTeX users please use
%\begin{thebibliography}{}
%
% and use \bibitem to create references. Consult the Instructions
% for authors for reference list style.
%
%\bibitem{RefJ}
% Format for Journal Reference
%Author, Article title, Journal, Volume, page numbers (year)
% Format for books
%\bibitem{RefB}
%Author, Book title, page numbers. Publisher, place (year)
% etc
%\end{thebibliography}

\end{document}
% end of file template.tex

